\section*{第 \ref{ch:b2:global-change} 章——全球变化与生物圈}

\begin{solution}{exo:b2:global-change:1}
$\Delta T = \lambda C_{\text{累积}}$，其中 $\lambda = \qty{1.65}{\celsius}$
每 \qty{1000}{GtC}：\qty{1.2}{\celsius}，\qty{1.65}{\celsius}，
\qty{3.3}{\celsius}。
\end{solution}

\begin{solution}{exo:b2:global-change:2}
度日是指：对各日的日均温超过某一基温的部分逐日求和，
即支配植物与变温动物发育的\omterm{thm:b2:global-change:phenology}{热时间}。
错配是指：当两个相互作用的物种依据不同信号安排季节，
而气候变暖只推动其中一方——以日长为信号的捕蝇鸟在非洲，
仍按旧日期抵达，
却发现以温度为信号的毛毛虫高峰
已经过去两周。
\end{solution}

\begin{solution}{exo:b2:global-change:3}
每升温一度，等温线约上升 \qty{150}{m}、向极地推移 \qty{150}{km}。
位于山顶的物种无更高之地可退：其等温线升至山顶以上，
栖息地随之消失，
而各温度带的面积随海拔升高而逐渐缩小。
\end{solution}

\begin{solution}{exo:b2:global-change:4}
溶解的 \ensuremath{\mathrm{CO_2}} 形成碳酸，碳酸
释放氢离子，pH 随 \ensuremath{\mathrm{CO_2}} 分压的
对数成比例下降。额外的氢离子使碳酸根转化为碳酸氢根，
因而碳酸根浓度
及碳酸钙的\omterm{thm:b2:global-change:acid}{饱和度}下降。
能构建文石的生物——珊瑚、翼足类、许多幼虫——
以及构建方解石的生物——颗石藻、有孔虫、软体动物——
构建壳体需付出更多代价，在\omterm{thm:b2:global-change:acid}{饱和度}以下则会丧失壳体。
\end{solution}

\begin{solution}{exo:b2:global-change:5}
\qty{1.5}{\celsius}：$1.5/1.65\times 1000 = \qty{909}{GtC}$，共计
尚余 \qty{209}{GtC}，可维持 21 年。\qty{2}{\celsius}：
\qty{1212}{GtC}，尚余 \qty{512}{GtC}，可维持 51 年。
\end{solution}

\begin{solution}{exo:b2:global-change:6}
$\sqrt{2\times 150/0.12} = 50$ 天；提前 $1.5/0.12 = 12.5$
天。
\end{solution}

\begin{solution}{exo:b2:global-change:7}
\qty{3}{\celsius}：等温线上升 $3000/6.5 = \qty{460}{m}$；
温度带移至 \qtyrange{1260}{1860}{m}，但山峰止于 \qty{1800}{m}：
温度带顶部的 \qty{60}{m} 将失去，其余部分所在处的面积
为原来的 $2^{-460/300} = 0.35$ 倍——约损失了三分之二的面积。
\qty{5}{\celsius}：\qty{770}{m}；温度带将为 \qtyrange{1570}{2170}{m}；
仅 \qtyrange{1570}{1800}{m} 存在，宽度为温度带的三分之一，
面积为原来的 $2^{-770/300} = 0.17$ 倍——
即不足原面积的十分之一：该物种实际上已经
消失。
\end{solution}

\begin{solution}{exo:b2:global-change:8}
560：pH $7.93$，氢离子 $\times 1.9$；800：$7.79$，$\times
2.6$；1000：$7.70$，$\times 3.1$。
\end{solution}

\begin{solution}{exo:b2:global-change:9}
$1 - (1 - h)^{0.25}$：\qty{26}{\%}、\qty{53}{\%}、\qty{82}{\%}。
短期损失较小，因为残存生境仍维持着
多数物种的种群；这些种群过小而难以存续，
数十年内相继灭绝——即\omterm{thm:b2:global-change:sar}{灭绝债务}。
\end{solution}

\begin{solution}{exo:b2:global-change:10}
$c = 2\sqrt{0.7\times 20} = \qty{7.5}{km/yr}$；\qty{1000}{km}
约需 130 年。将 $r$ 或 $D$ 减半，都会使 $c$ 除以
$\sqrt 2$，即降至 \qty{5.3}{km/yr}；两者都无法使其减半——
要使速度减半，须使乘积 $rD$ 减为四分之一。
\end{solution}

\begin{solution}{exo:b2:global-change:11}
本底：$8\times 10^{6}\times 10^{-6} = 8$ 种/年，
即每世纪 800 种；若为本底的一千倍，每年 8000 种，每世纪 \num{800000} 种。
一世纪内记录在案的一万种约为本底的十倍，
而非一千倍——但记录所覆盖的只是
已被描述并受监测的物种中的百分之几（鸟类、
哺乳动物、少数植物），其中速率确实达到本底的一百倍
或更高；对于尚未描述的大多数物种，多数
灭绝从未被观察到，只能借助物种–面积估算，
它是唯一可用的手段。真实情况介于记录数与估算值之间，
这正是所引用范围如此之宽的原因。
\end{solution}

\begin{solution}{exo:b2:global-change:12}
温度与\omterm{prop:b2:global-change:tcre}{累积排放}量成正比；
净排放为零时，累积总量停止增长，
温度也随之停止上升。在\omterm{thm:b2:biogeochemical-cycles:box}{箱式模型}中，排放为零时，
大气中超出的部分随汇的作用开始减少，
这本应导致降温——但海洋仍在趋向与升高的 \ensuremath{\mathrm{CO_2}} 平衡，
仍使表层继续升温，两种效应几乎相互抵消：温度
数个世纪内保持不变。论证中缺失的还有其他温室气体
（甲烷寿命较短，排放停止后其增温作用迅速减弱），
目前掩盖部分增温的气溶胶
（在产生它们的排放停止后数周内即消失），
以及反馈过程——永久冻土碳释放、森林衰亡——
它们自身还可能带来
额外的排放。
\end{solution}

\begin{solution}{pb:b2:global-change:1}
\textbf{1.} A：$700 + 80\times 10 = \qty{1500}{GtC}$。B：$700 +
\tfrac12\times 50\times 10 = \qty{950}{GtC}$。
\textbf{2.} A：$1.65\times 1.5 = \qty{2.5}{\celsius}$。B：
$1.65\times 0.95 = \qty{1.6}{\celsius}$。
\textbf{3.} 2050。A：\qty{1000}{GtC}，\qty{1.65}{\celsius}。B：
排放量 $10\int_0^{30}(1 - t/50)\,\mathrm{d}t = 10(30 - 9) =
\qty{210}{GtC}$，总计 910，\qty{1.5}{\celsius}。
\textbf{4.} $\lambda\times 100 = \qty{0.165}{\celsius}$，每十年
略低于观测的 $0.2$（观测值还包括其他温室气体
及气溶胶遮蔽作用的消失）。
\textbf{5.} 是的：\omterm{prop:b2:global-change:tcre}{累积排放}量在 2070 年停止增长，
温度也随之停止上升，仍为 \qty{1.6}{\celsius}——这
表明温度由总量决定，
而非由排放速率决定。
\textbf{6.} $2/1.65\times 1000 = \qty{1212}{GtC}$；路径 A 达到它
需 $512/10 = 51$ 年，于 2071 年。
\textbf{7.} $\sqrt{2\times 200/0.1} = 63$ 天。
\textbf{8.} 相对 2020 年的升温：A $2.5 - 1.2 =
\qty{1.3}{\celsius}$，B \qty{0.4}{\celsius}。提前量 $\Delta T/b$：
A 13 天，B 4 天。
\textbf{9.} 间隔：A $20 + 13 = 33$ 天；B 24 天。
\textbf{10.} 毛毛虫在展叶后第 15 至第 25 天
可供取食。B：鸟在第 24 天抵达，恰好在时段之内。
A：第 33 天抵达，比最后一条毛毛虫晚 8 天。
\textbf{11.} 鸟提前 $3\times 1.3 = 4$ 天：间隔 $33 - 4 =
29$ 天——仍在时段之外。
\textbf{12.} 这种鸟能够适应新的温度；它无法做到的是
在食物已经消失时喂养雏鸟，因为它的信号
与猎物的信号对变暖的反应并不相同。
\textbf{13.} A：$1.3/6.5\times 1000 = \qty{200}{m}$ 上移，
$1.3/0.65\times 100 = \qty{200}{km}$ 向极地推移。B：\qty{57}{m} 与
\qty{57}{km}。
\textbf{14.} 80 年间，树木仅移动 \qty{8}{km}：它无法
跟上 \qty{200}{km}（A），也达不到 \qty{57}{km}（B）——
树木存在滞后，两条路径上其分布区都将与气候
长期失衡，历经数个世纪。
\textbf{15.} A：温度带将为 \qtyrange{1700}{2000}{m}：恰好
到达山顶，已无更远之处可去。山体随
海拔升高而收窄，同样一段温度带只覆盖其原面积的
$2^{-200/300} = 0.63$ 倍——
此后任何进一步增温都将无地可供其迁入。
\textbf{16.} B：上移 \qty{57}{m}，面积因子 $2^{-57/300} = 0.88$：
即损失 \qty{12}{\%}。
\textbf{17.} 八十年间，以每十年 \qty{30}{km} 计：\qty{240}{km}，超过
A 路径所需的 \qty{200}{km}：该鱼类能跟得上（并离开了
其原有的渔场）。
\textbf{18.} 在较低的边缘，该物种遇到从下方迁来的竞争者、
病原体和捕食者，这些是它此前未曾面对的；
其自身适应旧温度的生理特性
被后者超越：这种退缩的竞争成分
与生理成分同样重要。
\textbf{19.} A：$\mathrm{pH} = 8.2 - 0.9\log_{10}(600/280) = 7.90$；
氢离子为 1750 年时的 $10^{0.30} = 2.0$ 倍。B：$8.04$；$1.44$
倍。
\textbf{20.} $\Omega = 4/(\text{比值})$：A $2.0$，B $2.8$。在 2 时
珊瑚仍能钙化，但速度缓慢，骨骼较弱，并且
礁体侵蚀速度超过生长速度；在 $2.8$ 时，珊瑚能够造礁，
不过路径 B 的高温仍会使其白化。
\textbf{21.} A：$400\times 0.4^{0.25} = 318$，损失 82 种。B：
$400\times 0.8^{0.25} = 378$，损失 22 种。
\textbf{22.} A：$318\times(1 - 0.05\times 1.3) = 298$。B：
$378\times(1 - 0.05\times 0.4) = 371$。
\textbf{23.} A：$1 - 298/400 = \qty{26}{\%}$。B：\qty{7}{\%}。
\textbf{24.} 两条路径均为生境清除：A 为 82 种对 20 种，B 为 22 种对
7 种——模型中的增温项影响较小；实际上
两者相互作用，因为破碎化的森林无法迁移其分布区。
\textbf{25.} 路径 A：\qty{2.5}{\celsius}，pH $7.90$，树种的四分之一
已经丧失。路径 B：\qty{1.6}{\celsius}，pH $8.04$，\qty{7}{\%}
的树种损失。
\end{solution}
